\documentclass[final]{beamer}

% ====================
% Packages
% ====================

\usepackage[T1]{fontenc}
\usepackage{lmodern}
\usepackage[size=custom, width=122, height=91, scale=1.0]{beamerposter}
\usetheme{gemini}
\usepackage{graphicx}
\usepackage{booktabs}
\usepackage{tikz}
\usetikzlibrary{calc,shadows}
\usepackage{pgfplots}
\pgfplotsset{compat=1.14}
\usepackage{anyfontsize}
\usepackage{ragged2e}
\usepackage{xcolor}
\usepackage{enumitem}
\usepackage{caption}
\usepackage[most]{tcolorbox}

% ====================
% Colors
% ====================

\definecolor{accent}{RGB}{196,80,0}
\definecolor{accentdark}{RGB}{139,50,0}
\definecolor{textmain}{RGB}{26,18,16}
\definecolor{textmid}{RGB}{74,58,44}
\definecolor{textdim}{RGB}{138,122,108}
\definecolor{panelborder}{RGB}{221,208,192}
\definecolor{bgwarm}{RGB}{255,250,245}

% ====================
% Override gemini headline colors  (default: fg=lightgray bg=blue — invisible on light bg)
% ====================

\setbeamercolor{headline}{fg=textmain, bg=bgwarm}
\setbeamercolor{headline title}{fg=textmain}
\setbeamercolor{headline author}{fg=textmid}
\setbeamercolor{headline institute}{fg=accent}
\setbeamercolor{headline rule}{bg=accent}

% ====================
% Font overrides
% ====================

\setsansfont{Lato}[
  UprightFont=*-Regular,
  ItalicFont=*-Italic,
  BoldFont=*-Medium,
  BoldItalicFont=*-MediumItalic
]

% ====================
% Figure caption helper
% ====================

\newcommand{\figcap}[1]{%
  \vspace{3pt}%
  {\centering\scriptsize\itshape\color{textdim}#1\par}%
}

% ====================
% Research-question numbered circle
% ====================

\newcommand{\rqnum}[1]{%
  \tikz[baseline=(n.base)]{%
    \node[circle, fill=accent, text=white,
          font=\scriptsize\bfseries\sffamily,
          minimum size=1.4em, inner sep=0pt] (n) {#1};%
  }%
}

% ====================
% Section title
% ====================

\newcommand{\sectitle}[1]{%
  {\usebeamerfont{block title}\color{accentdark}\MakeUppercase{#1}}%
  \par\nointerlineskip\vspace{2pt}%
  {\color{accent}\rule{\linewidth}{1.5pt}}%
  \par\vspace{4pt}%
}

% ====================
% Glassy panel
% ====================

\tcbset{
  glassy/.style={
    enhanced,
    colback=bgwarm,
    opacityback=0.72,
    colframe=panelborder,
    boxrule=0.5pt,
    arc=4pt,
    left=7pt, right=7pt, top=5pt, bottom=6pt,
    drop shadow,
    before upper={\color{textmain}\setlength{\parskip}{4pt}},
  }
}

% ====================
% Column widths
% ====================

\newlength{\sepwidth}
\newlength{\leftcolwidth}
\newlength{\midcolwidth}
\newlength{\rightcolwidth}

\setlength{\sepwidth}{0.016\paperwidth}
\setlength{\leftcolwidth}{0.300\paperwidth}
\setlength{\midcolwidth}{0.338\paperwidth}
\setlength{\rightcolwidth}{0.294\paperwidth}

\newcommand{\separatorcolumn}{\begin{column}{\sepwidth}\end{column}}

% ====================
% Background
% ====================

\setbeamertemplate{background}{%
  \begin{tikzpicture}
    \useasboundingbox (0,0) rectangle (\paperwidth,\paperheight);
    \node[anchor=south west, inner sep=0] at (0,0)
      {\includegraphics[width=\paperwidth,height=\paperheight]{background.jpg}};
    \fill[bgwarm,opacity=0.84] (0,0) rectangle (\paperwidth,\paperheight);
  \end{tikzpicture}%
}

% ====================
% Compact Headline
% (avoids tikzpicture width issues by using beamercolorbox for the stripe)
% ====================

\setbeamertemplate{headline}{%
  % Top accent stripe — solid beamercolorbox, no tikzpicture
  \begin{beamercolorbox}[wd=\paperwidth, ht=0.28cm, dp=0pt]{headline rule}%
  \end{beamercolorbox}%
  % Header body
  \begin{beamercolorbox}[wd=\paperwidth, ht=6.2cm, dp=0.35cm]{headline}%
    \hspace{1.6cm}%
    % Left logos
    \begin{minipage}[c][6.2cm][c]{5.5cm}
      \centering
      
\includegraphics[height=4cm]{msulogo.png}
    \end{minipage}%
    \hfill
    % Center title block
    \begin{minipage}[c][6.2cm][c]{0.56\paperwidth}
      \centering
      {\usebeamerfont{headline title}\usebeamercolor[fg]{headline title}%
       Active Solar Irradiance Observer \textcolor{accent}{(ASIO)}}\\[0.3ex]
      {\fontsize{10}{13}\selectfont\color{textmid}%
       Evan Williams $\cdot$ Sarah Pawlowski $\cdot$ Chris Entzel $\cdot$
       Charles Kankelborg $\cdot$ Camrynn Wells $\cdot$ Caleb Pascual $\cdot$
       Reed Lando $\cdot$ Jake Davis $\cdot$ Chase Bronson $\cdot$
       Tyler Holliday $\cdot$ Emily Schwartz $\cdot$ Aidan Rasmussen $\cdot$
       Connor Parrott $\cdot$ Larry Springer $\cdot$ Ian Garcia $\cdot$
       Jacob Weber $\cdot$ Finnegan Davidson $\cdot$ Sophia Daly}\\[0.4ex]
      {\fontsize{8.5}{10.5}\selectfont\color{accent}%
       \addfontfeatures{LetterSpace=5}\MakeUppercase{%
       Montana State University \,---\, Department of Physics
       \,---\, Space Science \& Engineering Laboratory}}
    \end{minipage}%
    \hfill
    % Right logos
    \begin{minipage}[c][6.2cm][c]{8cm}
      \centering
      
\includegraphics[height=3.8cm]{nasalogo.png}\hspace{4pt}%
      
\includegraphics[height=2.8cm]{muselogo.png}\hspace{4pt}%
      
\includegraphics[height=2.8cm]{asiologo.png}
    \end{minipage}%
    \hspace{1.6cm}%
  \end{beamercolorbox}%
  % Bottom accent rule
  \begin{beamercolorbox}[wd=\paperwidth, ht=0pt, dp=3pt]{headline rule}%
  \end{beamercolorbox}%
}

% Remove footer
\setbeamertemplate{footline}{}

% ====================
% Document
% ====================

\begin{document}
\begin{frame}[t]
\vspace{0.4cm}
\begin{columns}[t, totalwidth=\paperwidth]
\separatorcolumn

% ============================================================
%  LEFT COLUMN
% ============================================================
\begin{column}{\leftcolwidth}

% -- Introduction --
\begin{tcolorbox}[glassy]
\sectitle{Introduction}

Solar flares are among the most energetic events in our solar system,
releasing enormous quantities of energy across the electromagnetic
spectrum in timescales ranging from milliseconds to hours. The underlying
physical mechanisms --- particularly sub-second oscillatory phenomena
known as \textbf{\textcolor{accentdark}{Quasi-Periodic Pulsations (QPPs)}}
--- remain poorly understood because existing solar monitors lack the
temporal resolution and multi-band simultaneity needed to resolve them.

QPPs have been detected in microwave, hard X-ray, and EUV emissions,
but the temporal relationships \textit{between} these bands at sub-second
cadence are largely unconstrained. The spatial dynamics of the soft X-ray
source during a flare --- and the thermal evolution of the flare late
phase on centi-second timescales --- are similarly unexplored
observational frontiers.

The \textbf{\textcolor{accentdark}{Active Solar Irradiance Observer (ASIO)}}
is designed specifically to fill these gaps. Scheduled for launch in 2027
aboard NASA's \textbf{\textcolor{accentdark}{MUSE (MUlti-slit Solar
Explorer)}} spacecraft, ASIO delivers simultaneous high-cadence irradiance
measurements in the SXR and HXR bands with integrated coarse source
localization --- a capability not available in any existing solar instrument.
\end{tcolorbox}

\vspace{0.3cm}

% -- Research Objectives --
\begin{tcolorbox}[glassy]
\sectitle{Research Objectives}

\vspace{2pt}
\noindent\rqnum{1}\hspace{5pt}%
\begin{minipage}[t]{0.88\linewidth}
  Do sub-second \textbf{\textcolor{accentdark}{Quasi-Periodic Pulsations
  (QPPs)}} exist simultaneously across the SXR, HXR, and EUV energy
  bands? Are there measurable temporal offsets between bands that reveal
  the causal sequence of energy release?
\end{minipage}

\vspace{5pt}
\noindent\rqnum{2}\hspace{5pt}%
\begin{minipage}[t]{0.88\linewidth}
  Are \textbf{\textcolor{accentdark}{soft X-ray sources}} during solar
  flares spatially stationary, or do they exhibit rapid motion across the
  solar disk? If dynamic, over what spatial extent and on what timescale?
\end{minipage}

\vspace{5pt}
\noindent\rqnum{3}\hspace{5pt}%
\begin{minipage}[t]{0.88\linewidth}
  How does the \textbf{\textcolor{accentdark}{thermal structure}} of a
  solar flare evolve during its late phase on centi-second timescales ---
  and what temperature and emission measure trajectories emerge as the
  plasma cools?
\end{minipage}

\vspace{7pt}
\noindent
\begin{minipage}[t]{0.49\linewidth}
  \centering
  \includegraphics[width=\linewidth, height=4.5cm,
    keepaspectratio=false, clip]{slide2_img1.jpg}
  \figcap{\textbf{Fig.~1.} EUV observation of a solar flare. Scale
    markers indicate the source spans 50--190~Mm across the solar disk.}
\end{minipage}\hfill%
\begin{minipage}[t]{0.49\linewidth}
  \centering
  \includegraphics[width=\linewidth, height=4.5cm,
    keepaspectratio]{slide2_img2.png}
  \figcap{\textbf{Fig.~2.} QPP classification by period and duration
    (Tan et al.\ 2010). ASIO targets the fast and slow VSP regimes at
    sub-second cadence.}
\end{minipage}
\end{tcolorbox}

\vspace{0.3cm}

% -- References --
\begin{tcolorbox}[glassy]
\sectitle{References}
\color{textmid}
\scriptsize
\begin{enumerate}[leftmargin=16pt, itemsep=3pt,
    label={\textcolor{accent}{\bfseries\arabic*.}}]
  \item Ning, H., Chen, Y., Wu, Z., Su, Y., Tian, H., Li, G., Du, G.,
    \& Song, H. (2018). Two-stage energy release process of a confined
    flare with double HXR peaks. \textit{The Astrophysical Journal},
    \textit{854}(2), 178.
  \item De~Pontieu, B., Cheung, M., \& MUSE team. (2022). The MUlti-slit
    Solar Explorer (MUSE). \textit{Bulletin of the AAS}, \textit{54}(7).
  \item Tan, B., Zhang, Y., Tan, C., \& Liu, Y. (2010). Microwave
    quasi-periodic pulsations in multi-timescales associated with a
    solar flare/CME event. \textit{The Astrophysical Journal},
    \textit{723}(1), 25--39.
\end{enumerate}
\end{tcolorbox}

\end{column}
\separatorcolumn

% ============================================================
%  MIDDLE COLUMN
% ============================================================
\begin{column}{\midcolwidth}

% -- The Instrument --
\begin{tcolorbox}[glassy]
\sectitle{The Instrument}

ASIO is a compact photometer engineered as a hosted payload on the
\textbf{\textcolor{accentdark}{NASA MUSE spacecraft}}, which will occupy
a Sun-synchronous low Earth orbit at $\sim$667~km. The instrument leverages
MUSE's spacecraft bus for power, thermal control, and telemetry while
providing a high-cadence X-ray measurement channel fully independent of
MUSE's primary UV spectrograph.

The flight chassis is precision-machined from aluminum with multiple
aperture channels, each coupling solar X-rays to dedicated solid-state
photodiode detectors. A \textbf{\textcolor{accentdark}{shadow plate}}
upstream of the SXR detector array provides coarse source localization
without imaging optics --- differential shadow patterns across detector
segments encode the angular position of the X-ray source on the solar
disk. An adjustable mounting ring permits fine alignment of the instrument
axis after spacecraft integration.

\vspace{5pt}
\noindent
\begin{minipage}[t]{0.535\linewidth}
  \centering
  \includegraphics[width=\linewidth, height=4.8cm,
    keepaspectratio=false, clip]{slide3_img3.jpg}
  \figcap{\textbf{Fig.~3.} ASIO flight instrument --- precision-machined
    aluminum housing with multiple X-ray aperture channels and solid-state
    photodiode detectors.}
\end{minipage}\hfill%
\begin{minipage}[t]{0.445\linewidth}
  \centering
  \includegraphics[width=\linewidth, height=4.8cm,
    keepaspectratio]{slide3_img4.png}
  \figcap{\textbf{Fig.~4.} MUSE spacecraft CAD rendering. ASIO flies as a
    hosted payload on the main bus, launching in~2027.}
\end{minipage}

\vspace{5pt}
\noindent
\begin{minipage}[t]{0.49\linewidth}
  \centering
  \includegraphics[width=\linewidth, height=4cm,
    keepaspectratio]{slide3_img5.png}
  \figcap{\textbf{Fig.~5.} ASIO structural CAD model showing aperture
    geometry and detector mounting.}
\end{minipage}\hfill%
\begin{minipage}[t]{0.49\linewidth}
  \centering
  \includegraphics[width=\linewidth, height=4cm,
    keepaspectratio]{slide3_img6.png}
  \figcap{\textbf{Fig.~6.} Adjustable positioning ring enabling precise
    instrument alignment after spacecraft integration.}
\end{minipage}
\end{tcolorbox}

\vspace{0.3cm}

% -- Alignment Procedure --
\begin{tcolorbox}[glassy]
\sectitle{Alignment Procedure}

\noindent
\begin{minipage}[t]{0.62\linewidth}
  The shadow plate upstream of the SXR photodiode array encodes the
  angular position of a solar X-ray source as a differential shadow
  pattern. Comparing flux ratios between detector segments yields a
  coarse source centroid without imaging optics.
  Post-integration alignment is characterised with an automated
  edge-detection pipeline: exterior and interior aperture corners are
  located to sub-pixel accuracy, compared against the optical model,
  and the adjustable ring is locked once verified.
\end{minipage}\hfill%
\begin{minipage}[t]{0.35\linewidth}
  \centering
  \includegraphics[width=\linewidth, height=4.8cm,
    keepaspectratio]{slide5_img10.png}
  \figcap{\textbf{Fig.~7.} Edge-detection alignment output.}
\end{minipage}
\end{tcolorbox}

\end{column}
\separatorcolumn

% ============================================================
%  RIGHT COLUMN
% ============================================================
\begin{column}{\rightcolwidth}

% -- Instrument Testing --
\begin{tcolorbox}[glassy]
\sectitle{Instrument Testing}

All ASIO assembly and integration was performed under
\textbf{\textcolor{accentdark}{Class~100 cleanroom conditions}} at MSU's
SSEL facility. Calibrated red laser sources illuminated the aperture
assembly on an optical bench to map the spatial response of each
photodiode channel and validate the shadow-plate positional sensing
concept. Differential illumination patterns matched ray-trace model
predictions, confirming readiness for flight integration.

\vspace{5pt}
\noindent
\begin{minipage}[t]{0.32\linewidth}
  \centering
  \includegraphics[width=\linewidth, height=3.8cm,
    keepaspectratio=false, clip]{slide4_img7.jpg}
  \figcap{\textbf{Fig.~8.} ASIO integration in the Class~100 cleanroom
    facility at MSU~SSEL.}
\end{minipage}\hfill%
\begin{minipage}[t]{0.32\linewidth}
  \centering
  \includegraphics[width=\linewidth, height=3.8cm,
    keepaspectratio=false, clip, angle=270]{slide4_img9.jpg}
  \figcap{\textbf{Fig.~9.} ASIO on the optical bench during laser aperture
    characterization.}
\end{minipage}\hfill%
\begin{minipage}[t]{0.32\linewidth}
  \centering
  \includegraphics[width=\linewidth, height=3.8cm,
    keepaspectratio=false, clip]{slide4_img8.jpg}
  \figcap{\textbf{Fig.~10.} Red laser illumination of the shadow-plate
    aperture validating positional sensing geometry.}
\end{minipage}
\end{tcolorbox}

\vspace{0.3cm}

% -- Data Analysis --
\begin{tcolorbox}[glassy]
\sectitle{Data Analysis}

ASIO's science pipeline converts multi-band irradiance time series into
\textbf{\textcolor{accentdark}{temperature}} and
\textbf{\textcolor{accentdark}{emission measure}} using ratios of SXR and
HXR fluxes under an isothermal plasma assumption. Quasi-Periodic Pulsations
are identified via \textbf{\textcolor{accentdark}{continuous wavelet
decomposition}}, which resolves the time-frequency structure of the
irradiance signal and isolates statistically significant periodic components
above global (95\%) and local (99\%) significance thresholds. Synthetic
signal validation confirms robust QPP recovery from realistic noise
environments prior to first-light solar observations.

\vspace{5pt}
\noindent
\begin{minipage}[t]{0.49\linewidth}
  \centering
  \includegraphics[width=\linewidth, height=3cm,
    keepaspectratio]{slide6_img11.png}
  \figcap{\textbf{Fig.~11.} Simulated irradiance time series with embedded
    quasi-periodic structure and additive noise.}
\end{minipage}\hfill%
\begin{minipage}[t]{0.49\linewidth}
  \centering
  \includegraphics[width=\linewidth, height=3cm,
    keepaspectratio=false, clip]{slide6_img12.png}
  \figcap{\textbf{Fig.~12.} Rectified wavelet power spectrum. QPP periods
    recovered above the 95\% significance contour (red).}
\end{minipage}

\vspace{4pt}
\centering
\includegraphics[width=0.96\linewidth, height=3.8cm,
  keepaspectratio]{slide6_img13.png}
\figcap{\textbf{Fig.~13.} Reconstruction from significant wavelet
  coefficients only (RMSE $= 1.86\times10^{-2}$, 6.7\% of std),
  demonstrating high-fidelity QPP isolation from background noise.}
\end{tcolorbox}

\vspace{0.3cm}

% -- Acknowledgments --
\begin{tcolorbox}[glassy]
\sectitle{Acknowledgments}

This work is supported by NASA under the MUSE Mission development program.
Additional support provided by the Montana State University Space Science
\& Engineering Laboratory (SSEL) and the MSU Department of Physics.
Presented at the \textbf{MSU Student Research Celebration --- Spring 2026}.
\end{tcolorbox}

\end{column}
\separatorcolumn
\end{columns}
\end{frame}
\end{document}
